% sections/how_it_works.tex
% Seção 2: Como o Sistema Funciona

\section{Como o Sistema Funciona}

O IntelliStock opera através de quatro módulos de análise integrados,
chamados \textbf{Esferas}. Cada esfera tem uma responsabilidade
específica e gera alertas direcionados.

\subsection{Arquitetura de Quatro Esferas}

\begin{center}
\begin{tabular}{>{\bfseries}l p{8cm}}
\toprule
Esfera & Responsabilidade \\
\midrule
1 - Produtos Prontos   & Analisa itens comprados e revendidos diretamente \\
2 - Ingredientes       & Explode receitas e analisa necessidade de insumos \\
3 - Produção           & Identifica necessidade de produção de semi-acabados \\
4 - Perecibilidade     & Monitora riscos de vencimento e desperdício \\
\bottomrule
\end{tabular}
\end{center}

\subsection{Esfera 1: Produtos Prontos}

Analisa itens comprados e vendidos sem transformação significativa.
Exemplos: bebidas, sobremesas industrializadas, produtos embalados.

\textbf{Funcionalidades:}

\begin{itemize}[leftmargin=2em]
  \item Cálculo de demanda média com base em histórico de vendas
  \item Determinação do ponto de reposição considerando lead time
  \item Cálculo de estoque de segurança baseado em variabilidade
  \item Geração de alertas de \textbf{compra urgente} ou \textbf{planejamento}
\end{itemize}

\textbf{Exemplo:}

\begin{quote}
\textit{Refrigerante X: demanda média de 50 un/dia, lead time de 2 dias,
estoque atual de 80 unidades. Ponto de reposição = 100 + margem.
Como 80 < 100, o sistema gera alerta de compra.}
\end{quote}

\subsection{Esfera 2: Ingredientes}

Gerencia a explosão de receitas (BOM - Bill of Materials).
Quando um prato é vendido, seus ingredientes são consumidos proporcionalmente.

\textbf{Funcionalidades:}

\begin{itemize}[leftmargin=2em]
  \item Identificação de ingredientes por prato
  \item Cálculo de demanda consolidada de cada ingrediente
  \item Geração de alertas de compra para insumos
\end{itemize}

\textbf{Exemplo:}

\begin{quote}
\textit{Filé à Parmegiana usa 300g de filé mignon. Demanda prevista: 20 pratos/dia.
Demanda de filé mignon: 6kg/dia. O sistema verifica estoque e lead time
para gerar o alerta apropriado.}
\end{quote}

\subsection{Esfera 3: Produção}

Analisa itens semi-acabados produzidos internamente antes do uso.
Exemplos: molhos, massas frescas, pré-preparos.

\textbf{Funcionalidades:}

\begin{itemize}[leftmargin=2em]
  \item Identificação de semi-acabados que precisam ser produzidos
  \item Diferenciação entre produção \textbf{reativa} (vendas do dia) e
        \textbf{planejada} (demanda futura)
  \item Priorização de alertas reativos sobre planejados
  \item Direcionamento de alertas para a cozinha
\end{itemize}

\textbf{Exemplo:}

\begin{quote}
\textit{Molho bolonhesa é semi-acabado. Houve consumo por 5 vendas de lasanha
hoje. Estoque baixo. Sistema gera alerta urgente para cozinha:
"Produzir molho bolonhesa - demanda reativa".}
\end{quote}

\subsection{Esfera 4: Perecibilidade Inteligente}

Camada de supervisão que analisa risco de desperdício considerando
validade e demanda prevista.

\textbf{Funcionalidades:}

\begin{itemize}[leftmargin=2em]
  \item Identificação de lotes próximos do vencimento
  \item Simulação de consumo usando lógica FEFO (First Expired, First Out)
  \item Cálculo de risco de desperdício
  \item Geração de alertas de gestão sobre riscos
\end{itemize}

\begin{infobox}[Regra de Precedência]
A Esfera 4 \textbf{informa} sobre riscos, mas não bloqueia alertas de compra.
A operação mantém autonomia nas decisões.
\end{infobox}

\textbf{Exemplo:}

\begin{quote}
\textit{10kg de queijo mussarela com validade em 3 dias. Demanda média: 2kg/dia.
Consumo previsto: 6kg. Risco de desperdício: 4kg. Alerta para gestão:
"Risco de perda de 4kg de mussarela".}
\end{quote}

\subsection{Níveis de Prioridade}

O sistema classifica alertas em três níveis:

\begin{center}
\begin{tabular}{>{\bfseries}l l p{7cm}}
\toprule
Prioridade & Indicação & Significado \\
\midrule
URGENT & Vermelho & Ação imediata necessária \\
PLAN   & Amarelo  & Atenção para os próximos dias \\
INFO   & Azul     & Informação para acompanhamento \\
\bottomrule
\end{tabular}
\end{center}

\subsection{Direcionamento por Responsável}

Cada alerta é direcionado ao responsável apropriado:

\begin{center}
\begin{tabular}{>{\bfseries}l p{9cm}}
\toprule
Responsável & Escopo de Alertas \\
\midrule
Compras   & Compra de produtos e ingredientes \\
Cozinha   & Produção de semi-acabados \\
Gestão    & Perecibilidade, produtos parados, visão estratégica \\
\bottomrule
\end{tabular}
\end{center}

\subsection{Dados Essenciais para Operação}

Duas informações são fundamentais para o funcionamento do IntelliStock:

\textbf{1. Histórico de Vendas}

O sistema utiliza dados de vendas para calcular demanda e identificar padrões.
Recomendação: mínimo de 30 dias de histórico.

\textbf{2. Lead Time de Fornecedores}

O tempo de entrega de cada fornecedor determina a antecedência necessária
para compras. Quanto menor o lead time, menor o estoque necessário.

\begin{infobox}[Configuração Inicial]
A qualidade dos alertas depende diretamente da qualidade dos dados.
Lead time e histórico de vendas são configurações essenciais.
\end{infobox}
